% ============================================================
% Capítulo 2 — Contexto del problema
% ============================================================

\chapter{Contexto del problema}
\label{cap:contexto}

El objetivo central del trabajo es estudiar qué aporta la topología a la
segmentación de estructuras finas: cómo se miden las formas, cómo se usan
como funciones de pérdida y qué mejoran frente a las pérdidas de
solapamiento. Para ponerlas a prueba necesitamos un caso práctico con
estructuras finas y errores topológicos claros. Este capítulo presenta ese
caso: la detección de filamentos solares. Describimos primero qué son los
filamentos y por qué interesa detectarlos (Sección~\ref{sec:filamentos}),
después el conjunto de datos MAGFiLO (Sección~\ref{sec:magfilo}) y la
competición en la que se enmarca el trabajo, con su métrica de evaluación
(Sección~\ref{sec:competicion}), y revisamos los trabajos previos de
detección automática (Sección~\ref{sec:estado-arte}). Terminamos con las
dificultades específicas del problema (Sección~\ref{sec:dificultades}), que
son en buena parte topológicas.

% ------------------------------------------------------------
\section{Filamentos solares}
\label{sec:filamentos}
% ------------------------------------------------------------

Los filamentos son nubes de plasma relativamente frío y denso suspendidas
en la corona solar por el campo magnético, sobre las líneas de inversión
de polaridad de la fotosfera~\cite{martin1998,mackay2010}. Vistos sobre el
disco solar en la línea H-alfa ($656{,}3$~nm), absorben la radiación del
fondo y aparecen como estructuras oscuras y alargadas. Cuando se observan
en el limbo, sobre el fondo oscuro del espacio, aparecen brillantes y
reciben el nombre de \emph{protuberancias}~\cite{parenti2014}.

Morfológicamente, un filamento consta de un cuerpo principal alargado, la
\emph{espina}, del que salen estructuras finas laterales llamadas
\emph{barbas}. La orientación de las barbas respecto a la espina permite
inferir la \emph{quiralidad} del filamento (dextrógira o levógira), que
refleja la configuración de su campo magnético~\cite{martin1998}.

Su interés práctico es el clima espacial. Los filamentos pueden
desestabilizarse y entrar en erupción, y están en el origen de eyecciones
de masa coronal, fulguraciones y tormentas de partículas
energéticas~\cite{schwenn2006}. Una eyección dirigida hacia la Tierra
puede dañar la red eléctrica, alterar los sistemas GPS y suponer un riesgo
de radiación para tripulaciones y astronautas~\cite{kaggle2026}. Detectar
y caracterizar los filamentos de forma automática y precisa es, por tanto,
un paso previo para su seguimiento y para la predicción de erupciones.

\paragraph{Observaciones GONG.} Las imágenes proceden de la red GONG
(\emph{Global Oscillation Network Group}), un conjunto de seis
observatorios terrestres repartidos en longitud para observar el Sol de
forma casi continua~\cite{harvey1996}: Big Bear (B), Cerro Tololo (C),
Learmonth (L), Mauna Loa (M), El Teide (T) y Udaipur (U). La letra entre
paréntesis es la que identifica la estación en el nombre de cada fichero.
Las imágenes H-alfa de disco completo tienen $2048\times2048$
píxeles~\cite{kaggle2026}. Al ser observaciones terrestres, están
afectadas por la turbulencia atmosférica, las nubes y artefactos
instrumentales, y su calidad varía entre estaciones y días.

% TODO: figura con una imagen H-alfa completa y sus anotaciones.
% \begin{figure}[ht]
%   \centering
%   \includegraphics[width=0.6\textwidth]{figuras/halpha_ejemplo.png}
%   \caption{Observación H-alfa de GONG con los filamentos anotados.}
%   \label{fig:halpha}
% \end{figure}

% ------------------------------------------------------------
\section{El conjunto de datos MAGFiLO}
\label{sec:magfilo}
% ------------------------------------------------------------

MAGFiLO (\emph{Manually Annotated GONG Filaments in H-alpha
Observations}) es un conjunto de datos de filamentos anotados
manualmente sobre observaciones H-alfa de GONG~\cite{ahmadzadeh2024}.
Su versión 1.0 contiene $10\,244$ filamentos anotados en $1\,593$
observaciones entre 2011 y 2022. Para cada filamento incluye la máscara
de segmentación, la caja contenedora mínima, la espina y la quiralidad.

La anotación supuso más de $1\,000$ horas-persona, con un proceso de
revisión doble ciego. Cada observación fue anotada por hasta tres
anotadores independientes y revisada después por expertos. El acuerdo
entre anotadores, medido con el estadístico $\kappa$, es de
$0{,}66$~\cite{ahmadzadeh2024}. Es un acuerdo sustancial pero lejos del
perfecto: la propia referencia humana es ruidosa, y esto condiciona
tanto el entrenamiento como la interpretación de las métricas.

\paragraph{Subconjunto de la competición.} La competición usa solo las
máscaras de segmentación~\cite{kaggle2026}. El conjunto de entrenamiento
público está organizado por \emph{pares anotador--imagen}: cada imagen
aparece tantas veces como anotadores la hayan etiquetado, cada vez con su
propio conjunto de filamentos. La Tabla~\ref{tab:datos} resume sus
características.

\begin{table}[ht]
\centering
\begin{tabular}{@{}lr@{}}
\toprule
Imágenes distintas & 707\\
Pares anotador--imagen & 1\,154\\
\quad imágenes con 1 / 2 / 3 anotadores & 411 / 145 / 151\\
Filamentos anotados (sumando anotadores) & 8\,199\\
Filamentos por par (media) & 7{,}1\\
Periodo & 2011--2022\\
Resolución & $2048\times2048$\\
\bottomrule
\end{tabular}
\caption{Conjunto de entrenamiento público de la competición, calculado a
partir del fichero de anotaciones.}
\label{tab:datos}
\end{table}

Las seis estaciones están representadas de forma equilibrada, entre 102 y
127 imágenes cada una. La distribución temporal no lo está: la mayoría de
las imágenes son de 2011--2017, y los años 2018--2022 aportan pocas
observaciones. Además, de un mismo día suele haber varias observaciones,
tomadas desde estaciones distintas, que muestran prácticamente el mismo
Sol. Este hecho obliga a separar entrenamiento y validación por fecha y
no por imagen, como se explica en el capítulo de metodología.

El conjunto de test es oculto. Solo se dispone de sus imágenes, y las
predicciones se evalúan en la plataforma.

% ------------------------------------------------------------
\section{La competición \emph{Solar Filament Segmentation Challenge 2026}}
\label{sec:competicion}
% ------------------------------------------------------------

El trabajo se enmarca en el \emph{Solar Filament Segmentation Challenge
2026}, una competición de Kaggle organizada dentro de la IEEE BigData Cup
2026 y patrocinada por el National Solar Observatory~\cite{kaggle2026}.
Se lanzó el 10 de julio de 2026 y el plazo de entrega final es el 15 de
noviembre de 2026.

\paragraph{Tarea.} Se pide generar una máscara de segmentación para cada
filamento \emph{individual}, que capture toda su extensión, incluidas las
estructuras finas como las barbas, y que incluya el menor número posible
de píxeles ajenos al filamento~\cite{kaggle2026}. Se trata, por tanto, de
un problema de \emph{segmentación de instancias} de una sola clase. El
envío es un fichero con una fila por filamento predicho, que contiene el
identificador de la imagen y la máscara codificada en RLE a resolución
original.

\paragraph{Métrica.} La métrica principal es la \emph{calidad
panóptica} (\emph{Panoptic Quality}, PQ)~\cite{kirillov2019}. Inicialmente
la competición usaba el coeficiente Dice, pero los organizadores lo
sustituyeron por PQ el 7 de agosto de 2026~\cite{kaggle2026}.

\begin{definicion}[{\cite{kirillov2019}}]
Sean $Y$ el conjunto de segmentos de referencia e $\hat Y$ el de segmentos
predichos. Un par $(y,\hat y)\in Y\times\hat Y$ es un \emph{verdadero
positivo} si $\operatorname{IoU}(y,\hat y)>0{,}5$, donde
\[
  \operatorname{IoU}(y,\hat y)=\frac{|y\cap\hat y|}{|y\cup\hat y|}.
\]
Los segmentos predichos sin pareja son \emph{falsos positivos} (FP) y los
de referencia sin pareja, \emph{falsos negativos} (FN). La calidad
panóptica es
\[
  \mathrm{PQ}=\frac{\sum_{(y,\hat y)\in \mathit{TP}}\operatorname{IoU}(y,\hat y)}
                   {|\mathit{TP}|+\tfrac12|\mathit{FP}|+\tfrac12|\mathit{FN}|}.
\]
\end{definicion}

El umbral $\operatorname{IoU}>0{,}5$ tiene una consecuencia importante. Si
los segmentos predichos no se solapan entre sí, cada segmento de
referencia puede emparejarse como mucho con uno predicho, y viceversa. El
emparejamiento es, por tanto, único y no requiere ningún algoritmo de
asignación~\cite{kirillov2019}. Además, PQ se factoriza como
\[
  \mathrm{PQ}=
  \underbrace{\frac{\sum_{\mathit{TP}}\operatorname{IoU}}{|\mathit{TP}|}}_{\mathrm{SQ}}
  \times
  \underbrace{\frac{|\mathit{TP}|}{|\mathit{TP}|+\tfrac12|\mathit{FP}|+\tfrac12|\mathit{FN}|}}_{\mathrm{RQ}},
\]
donde la \emph{calidad de segmentación} SQ mide lo bien delimitados que
están los aciertos y la \emph{calidad de reconocimiento} RQ es una
medida $F_1$ de detección~\cite{kirillov2019}.

En la competición, $Y$ e $\hat Y$ contienen los segmentos de
\emph{todas} las imágenes. Es decir, se suman TP, FP, FN e IoU sobre todo
el conjunto antes de dividir, y no se promedia imagen a
imagen~\cite{kaggle2026}. Como unidad de evaluación se usa el par
anotador--imagen: las predicciones de una imagen se comparan por separado
con la anotación de cada anotador.

Como referencia de escala, los organizadores indicaron que cualquier
puntuación PQ superior a $0{,}35$ les resulta de gran valor~\cite{kaggle2026}.

\paragraph{Rúbrica.} La clasificación en el \emph{leaderboard} es solo una
parte de la evaluación final~\cite{kaggle2026}:
\begin{itemize}
  \item \textbf{Comparación cuantitativa (70\,\%):} PQ y las
  distribuciones de Dice, de IoU y de relaciones uno-a-varios y
  varios-a-uno entre segmentos predichos y de referencia. Estas últimas
  penalizan la \emph{fragmentación} (un filamento predicho en varios
  trozos) y la \emph{fusión} (varios filamentos predichos como uno).
  \item \textbf{Comparación cualitativa (30\,\%):} la descripción del
  pipeline completo, la morfología aparente de las predicciones y la
  calidad del código.
\end{itemize}
El \emph{leaderboard} público se calcula con aproximadamente la mitad del
conjunto de test; la otra mitad solo la ven los organizadores.

% ------------------------------------------------------------
\section{Detección automática de filamentos: trabajos previos}
\label{sec:estado-arte}
% ------------------------------------------------------------

La detección automática de filamentos en imágenes H-alfa tiene una
trayectoria de más de dos décadas. Podemos distinguir dos etapas.

\paragraph{Métodos clásicos.} Los primeros métodos se basan en que los
filamentos son más oscuros que su entorno. Shih y Kowalski~\cite{shih2003}
extraen los filamentos con umbrales locales adaptativos y operaciones
morfológicas. Qu et al.~\cite{qu2005} combinan umbralización, crecimiento
de regiones y un clasificador SVM para descartar las manchas solares, que
también son oscuras. Bernasconi et al.~\cite{bernasconi2005} desarrollan
un sistema completo que detecta, sigue y caracteriza filamentos, incluida
su quiralidad a partir de las barbas. Hao et al.~\cite{hao2013} usan un
umbral adaptativo basado en el método de Otsu y aplican el método a un
ciclo solar completo de datos H-alfa. Estos métodos son interpretables y
rápidos, pero dependen de umbrales ajustados a mano, son sensibles a las
condiciones de observación y tienden a confundir filamentos con manchas
solares y otras regiones oscuras~\cite{diercke2024}.

\paragraph{Aprendizaje profundo.} Con las redes convolucionales, el
problema se ha abordado como segmentación semántica o de instancias.
Zhu et al.~\cite{zhu2019} entrenan una U-Net~\cite{ronneberger2015} sobre
imágenes H-alfa del observatorio de Big Bear con máscaras obtenidas de
forma semiautomática. Ahmadzadeh et al.~\cite{ahmadzadeh2019} usan Mask
R-CNN, que segmenta instancias, entrenada con máscaras derivadas de
detecciones automáticas previas en lugar de anotaciones manuales.
Guo et al.~\cite{guo2022} detectan y clasifican filamentos con una red de
segmentación de instancias. Diercke et al.~\cite{diercke2024} proponen un
método semisupervisado: un detector YOLOv5 entrenado con cajas anotadas a
mano genera etiquetas aproximadas sobre el archivo de GONG, y con ellas
entrenan una U-Net.

Un rasgo común de estos trabajos es la escasez de anotaciones manuales a
nivel de píxel: la mayoría entrena con máscaras obtenidas por
umbralización o por otros detectores, que heredan sus errores. MAGFiLO
cubre precisamente esta carencia~\cite{ahmadzadeh2024}. Sobre MAGFiLO,
Solomon et al.~\cite{solomon2025} proponen EdgeAttNet, una U-Net con
atención guiada por bordes en el cuello de botella, orientada a capturar
mejor las barbas, y la comparan con una U-Net estándar entrenada con
BCE + Dice.

Hasta donde conocemos, estos trabajos optimizan pérdidas de solapamiento
píxel a píxel y se evalúan con métricas de solapamiento (precisión,
exhaustividad, IoU). Ninguno incorpora restricciones topológicas al
entrenamiento ni al post-procesado. Tampoco se habían evaluado con una
métrica de instancias como PQ, que penaliza explícitamente la
fragmentación y la fusión. Este es el hueco que aborda este trabajo.

% ------------------------------------------------------------
\section{Dificultades del problema}
\label{sec:dificultades}
% ------------------------------------------------------------

Los propios organizadores señalan tres dificultades
principales~\cite{kaggle2026}, a las que añadimos dos más:

\begin{enumerate}
  \item \textbf{Estructuras finas.} Las barbas tienen pocos píxeles de
  anchura y contraste débil. Son las primeras en perderse, y su
  orientación es justo la que aporta información física.
  \item \textbf{Ruido e imagen terrestre.} Distinguir el material del
  filamento de otras regiones oscuras (manchas, ruido, artefactos) no es
  trivial en observaciones afectadas por la atmósfera.
  \item \textbf{Continuidad estructural.} Los algoritmos tienden a
  devolver filamentos \emph{fragmentados} o a agrupar regiones oscuras
  cercanas en «islas», con resultados físicamente incoherentes.
  \item \textbf{Desacuerdo entre anotadores.} Con $\kappa=0{,}66$, los
  anotadores no coinciden siempre en dónde empieza y acaba un filamento,
  ni en si dos trozos forman uno o dos filamentos. Esto fija un techo a lo
  que puede conseguir cualquier modelo frente a una anotación concreta.
  \item \textbf{Desequilibrio de clases.} Los filamentos ocupan una
  fracción muy pequeña de cada imagen, de modo que una pérdida píxel a
  píxel está dominada por el fondo.
\end{enumerate}

Las dificultades 1 y 3 son, en el fondo, \emph{topológicas}. Si las
instancias se obtienen como componentes conexas de una máscara binaria,
el número de filamentos predichos es exactamente el número de Betti
$\beta_0$ de esa máscara. Un filamento partido en dos aumenta $\beta_0$
y puede convertir un acierto en un falso negativo y dos falsos positivos,
aunque solo falten unos pocos píxeles. Dos filamentos unidos por un puente de píxeles lo
reducen y provocan el error contrario. Una barba perdida, en cambio, apenas
cambia el IoU. Métricas como Dice o el propio IoU, que cuentan píxeles,
son poco sensibles a estos errores, mientras que PQ y la rúbrica los
penalizan con fuerza. Esta discrepancia entre lo que optimiza una pérdida
de solapamiento y lo que mide la evaluación es la motivación para
aplicar a este caso las herramientas topológicas del Capítulo~\ref{cap:fundamentos}.
